\documentclass[12pt,a4paper]{article}
\usepackage{amsmath}
\usepackage{amssymb}
\usepackage{graphicx}
\usepackage{algorithm}
\usepackage{algpseudocode}
\usepackage{listings}
\usepackage{xcolor}
\usepackage{geometry}
\geometry{margin=1in}

\lstset{
    language=Python,
    basicstyle=\ttfamily\small,
    keywordstyle=\color{blue},
    commentstyle=\color{gray},
    stringstyle=\color{red},
    showstringspaces=false,
    breaklines=true,
    frame=single
}

\title{Statistical Shape Model (SSM) Deformation: \\
Technical Documentation}
\author{SSM Visualization Tool}
\date{\today}

\begin{document}

\maketitle

\section{Introduction}

This document explains the mathematical framework and implementation of Statistical Shape Model (SSM) deformation using Principal Component Analysis (PCA) with particle-based correspondence.

\section{Mathematical Framework}

\subsection{Shape Representation}

A shape is represented by $N$ particles in 3D space:
\begin{equation}
\mathbf{P} = \begin{bmatrix}
x_1 & y_1 & z_1 \\
x_2 & y_2 & z_2 \\
\vdots & \vdots & \vdots \\
x_N & y_N & z_N
\end{bmatrix} \in \mathbb{R}^{N \times 3}
\end{equation}

\subsection{PCA Decomposition}

The shape space is decomposed using PCA:
\begin{equation}
\mathbf{P} = \bar{\mathbf{P}} + \sum_{i=1}^{M} w_i \sqrt{\lambda_i} \mathbf{v}_i
\end{equation}

where:
\begin{itemize}
    \item $\bar{\mathbf{P}}$ is the mean shape (particles)
    \item $\lambda_i$ are eigenvalues (variance in mode $i$)
    \item $\mathbf{v}_i$ are eigenvectors (shape variation modes)
    \item $w_i$ are weights (typically $w_i \in [-3, 3]$ for $\pm 3\sigma$)
    \item $M$ is the number of modes
\end{itemize}

\subsection{Mesh Deformation via Interpolation}

Since particles are sparse, we interpolate deformation to dense mesh vertices using inverse distance weighting (IDW):

\begin{equation}
\mathbf{M}_{deformed} = \mathbf{M}_{mean} + \sum_{j=1}^{k} \alpha_j \cdot \Delta \mathbf{P}_j
\end{equation}

where:
\begin{itemize}
    \item $\mathbf{M}_{mean}$ is the mean mesh vertex
    \item $\Delta \mathbf{P}_j$ is displacement of $j$-th nearest particle
    \item $\alpha_j$ is the interpolation weight
    \item $k$ is the number of nearest neighbors (typically $k=3$)
\end{itemize}

The interpolation weights are computed as:
\begin{equation}
\alpha_j = \frac{w_j}{\sum_{l=1}^{k} w_l}, \quad w_j = \frac{1}{d_j + \epsilon}
\end{equation}

where $d_j$ is the distance to the $j$-th nearest particle and $\epsilon = 10^{-10}$ prevents division by zero.

\section{Algorithm}

\begin{algorithm}
\caption{SSM Shape Reconstruction}
\begin{algorithmic}[1]
\State \textbf{Input:} Weights $\mathbf{w} = [w_1, w_2, \ldots, w_M]$
\State \textbf{Output:} Deformed mesh $\mathbf{M}_{deformed}$

\State \textbf{Step 1: Deform Particles}
\State $\mathbf{P}_{deformed} \gets \mathbf{P}_{mean}$
\For{$i = 1$ to $M$}
    \State $\mathbf{P}_{deformed} \gets \mathbf{P}_{deformed} + w_i \sqrt{\lambda_i} \mathbf{v}_i$
\EndFor

\State \textbf{Step 2: Compute Particle Displacement}
\State $\Delta \mathbf{P} \gets \mathbf{P}_{deformed} - \mathbf{P}_{mean}$

\State \textbf{Step 3: Interpolate to Mesh Vertices}
\State $\mathbf{M}_{deformed} \gets \mathbf{M}_{mean}$
\For{each vertex $v$ in mesh}
    \State Find $k$ nearest particles to $v$
    \State Compute interpolation weights $\alpha_1, \ldots, \alpha_k$
    \State $\mathbf{M}_{deformed}[v] \gets \mathbf{M}_{mean}[v] + \sum_{j=1}^{k} \alpha_j \cdot \Delta \mathbf{P}_j$
\EndFor

\State \textbf{return} $\mathbf{M}_{deformed}$
\end{algorithmic}
\end{algorithm}

\section{Numerical Example}

\subsection{Setup}

Consider a simple 2D example with:
\begin{itemize}
    \item 3 particles: $\mathbf{P}_{mean} = \begin{bmatrix} 0 & 0 \\ 1 & 0 \\ 0.5 & 1 \end{bmatrix}$
    \item 1 mode with $\lambda_1 = 0.25$, $\mathbf{v}_1 = \begin{bmatrix} 0.1 & 0 \\ 0.1 & 0 \\ 0 & 0.2 \end{bmatrix}$
    \item Weight: $w_1 = 2.0$
\end{itemize}

\subsection{Step 1: Deform Particles}

\begin{align}
\mathbf{P}_{deformed} &= \mathbf{P}_{mean} + w_1 \sqrt{\lambda_1} \mathbf{v}_1 \\
&= \begin{bmatrix} 0 & 0 \\ 1 & 0 \\ 0.5 & 1 \end{bmatrix} + 2.0 \times \sqrt{0.25} \times \begin{bmatrix} 0.1 & 0 \\ 0.1 & 0 \\ 0 & 0.2 \end{bmatrix} \\
&= \begin{bmatrix} 0 & 0 \\ 1 & 0 \\ 0.5 & 1 \end{bmatrix} + \begin{bmatrix} 0.1 & 0 \\ 0.1 & 0 \\ 0 & 0.2 \end{bmatrix} \\
&= \begin{bmatrix} 0.1 & 0 \\ 1.1 & 0 \\ 0.5 & 1.2 \end{bmatrix}
\end{align}

\subsection{Step 2: Particle Displacement}

\begin{equation}
\Delta \mathbf{P} = \mathbf{P}_{deformed} - \mathbf{P}_{mean} = \begin{bmatrix} 0.1 & 0 \\ 0.1 & 0 \\ 0 & 0.2 \end{bmatrix}
\end{equation}

\subsection{Step 3: Interpolate to Mesh Vertex}

Consider a mesh vertex at $\mathbf{m} = [0.5, 0.3]$. Find 2 nearest particles:
\begin{itemize}
    \item Particle 1: distance $d_1 = 0.5$, displacement $\Delta \mathbf{P}_1 = [0.1, 0]$
    \item Particle 3: distance $d_2 = 0.7$, displacement $\Delta \mathbf{P}_3 = [0, 0.2]$
\end{itemize}

Compute weights:
\begin{align}
w_1 &= \frac{1}{0.5} = 2.0, \quad w_2 = \frac{1}{0.7} \approx 1.43 \\
\alpha_1 &= \frac{2.0}{2.0 + 1.43} \approx 0.583 \\
\alpha_2 &= \frac{1.43}{2.0 + 1.43} \approx 0.417
\end{align}

Interpolated displacement:
\begin{align}
\Delta \mathbf{m} &= \alpha_1 \Delta \mathbf{P}_1 + \alpha_2 \Delta \mathbf{P}_3 \\
&= 0.583 \times [0.1, 0] + 0.417 \times [0, 0.2] \\
&= [0.058, 0.083]
\end{align}

Final deformed vertex:
\begin{equation}
\mathbf{m}_{deformed} = [0.5, 0.3] + [0.058, 0.083] = [0.558, 0.383]
\end{equation}

\section{Implementation}

\subsection{Core Reconstruction Function}

\begin{lstlisting}
def reconstruct(self, weights):
    # Step 1: Deform particles
    particles_deformed = self.particles.copy()
    for i, w in enumerate(weights):
        particles_deformed += w * np.sqrt(self.eigenvalues[i]) * self.eigenvectors[i]
    
    # Step 2: Compute displacement
    particle_displacement = particles_deformed - self.particles
    mesh_deformed = self.mesh_vertices.copy()
    
    # Step 3: Interpolate to mesh
    for i in range(len(self.mesh_vertices)):
        displacement = np.sum(
            particle_displacement[self.nearest_particles[i]] * 
            self.interpolation_weights[i][:, np.newaxis],
            axis=0
        )
        mesh_deformed[i] += displacement
    
    return mesh_deformed
\end{lstlisting}

\subsection{Building Deformation Map}

\begin{lstlisting}
def _build_deformation_map(self):
    from scipy.spatial import cKDTree
    tree = cKDTree(self.particles)
    distances, indices = tree.query(self.mesh_vertices, k=3)
    
    # Inverse distance weighting
    weights = 1.0 / (distances + 1e-10)
    weights = weights / weights.sum(axis=1, keepdims=True)
    
    self.nearest_particles = indices
    self.interpolation_weights = weights
\end{lstlisting}

\section{Variance Explained}

Each mode captures a percentage of total variance:
\begin{equation}
\text{Variance}_i = \frac{\lambda_i}{\sum_{j=1}^{M} \lambda_j} \times 100\%
\end{equation}

Typically, the first few modes capture most variation:
\begin{itemize}
    \item Mode 0: $\sim$50\% (primary shape variation)
    \item Mode 1: $\sim$20\% (secondary variation)
    \item Mode 2: $\sim$10\% (tertiary variation)
    \item Remaining modes: $<$20\%
\end{itemize}

\section{Interactive Visualization}

The tool provides sliders for each mode with range $[-2, 2]$ representing $\pm 2\sigma$ (standard deviations). This covers approximately 95\% of natural shape variation.

\subsection{Performance Optimization}

\begin{itemize}
    \item \textbf{Render-on-release}: Updates only when slider is released, not during drag
    \item \textbf{Pre-computed mapping}: Nearest neighbors computed once during initialization
    \item \textbf{Vectorized operations}: NumPy arrays for efficient computation
\end{itemize}

\section{Conclusion}

The SSM deformation combines:
\begin{enumerate}
    \item PCA-based shape space representation
    \item Particle-based correspondence
    \item Inverse distance weighted interpolation
    \item Interactive real-time visualization
\end{enumerate}

This enables intuitive exploration of anatomical shape variations captured by the statistical model.

\end{document}
